\section{Funding: from a structural rate response to value}

Suppose treasury asks for the value of a ten-basis-point reduction in a
funding action over four quarters. Three objects are required. The consumer
must specify the action and exposure. A scientific analysis must establish
what change in paid funding rate that action would cause under its model.
A valuation step must then translate that response into dated cash flows.
A forecast of market rates alone supplies neither the intervention response
nor the exposure response to the bank's decision.

The inspected MacroFinance inventory directs active numerical work to
BayesFilter and describes the root migration adapter as a compatibility
facade~\citep{v12_in_inventory}. In the re-inspected BayesFilter declarations,
\texttt{BayesianSSMProblem} binds the model manifest, data signature, static
shape and coordinate convention; \texttt{FilterProgram} separately names the
filter, required model capabilities, target policy and approximation
semantics~\citep{exec13_i3_inspection}. A proposed funding adapter should retain
these distinctions. An exact structural model, a deterministic approximate
filter and a fixed-randomness numerical procedure need not imply the same
likelihood or uncertainty. Importing a historical local sampler because it
returns a familiar array could bypass the selected authority entirely.

For a transparent conditional example, fix exposure $B=100$ million USD,
action $a=-10$ basis points, annual effective discount rate $r=0.04$ and an
initial implementation cost $K=40{,}000$ USD. Let $\theta$ denote paid-rate
pass-through in basis points per action basis point, held constant across
four quarters. Assume balances are invariant to the action, cash accrues
uniformly within each quarter and is paid at quarter end, and no additional
fees or risk effects occur. The present value of savings is then
\begin{equation}
 V(\theta)=-B\frac{a}{10{,}000}\theta
      \sum_{q=1}^{4}\frac{1/4}{(1+r)^{q/4}}-K
       =97{,}584.369097\theta-40{,}000.
 \label{exec13:i3-funding}
\end{equation}
The sign is explicit: a negative rate change decreases cost and produces a
positive saving. The factor $10{,}000$ converts basis points to a decimal
rate, while $1/4$ converts an annualized rate into quarterly accrual. The
fractional discount exponent is required because $r$ is annual effective.
Replacing it with a quarterly rate would require a different parameterization.

For an assumed pass-through range $[0.6,0.9]$, the resulting value range is
$[18{,}550.62,47{,}825.93]$ USD; at $\theta=0.75$ the value is
$33{,}188.28$. The break-even coefficient is approximately $0.409902$.
These figures reproduce both by the affine formula and by summing the four
cash flows, including a dollars-to-millions round trip. They are generated
planning calculations with supplied inputs, not bank estimates or a
confidence interval~\citep{exec13_i3_examples}. The coefficient range could
instead be an identified set or a fitted confidence set, but only if a
scientific analysis justified and labeled that interpretation.

The example makes several adverse exchanges visible. Treating each quarter
as a full year of accrual changes the midpoint value to $252{,}753.11$.
Reading ten basis points as ten percentage points changes the rate scale by
a factor of one hundred. Reading millions as dollars changes the exposure
scale by a factor of one million. These are semantic errors even if the
downstream computation is numerically exact. A consumer acceptance fixture
should assert the declared dimensions and the independently computed cash
flows, not only compare two calls to the same valuation function.

More consequentially, an action may change the balance itself. With
action-dependent balance $B_q(a)$ and funding rate $s_q(a)$, the cost
contrast is based on
\[
 \sum_q d_q\Delta_q
       \{B_q(a)s_q(a)-B_q(0)s_q(0)\},
\]
where $d_q$ discounts the cash flow and $\Delta_q$ is the accrual fraction.
Equation~\eqref{exec13:i3-funding} follows only under its fixed-exposure
restriction and the specified rate contrast. Multiplying a rate effect by
the observed balance does not account for endogenous deposit attrition,
replacement funding or interaction between amount and price. A useful
adapter therefore either retains this restricted target or asks for the
joint structural response. It does not silently convert a restricted
scenario into a general treasury value claim.

The first proposed workflow is consequently modest and complete: a
MacroFinance question supplies target and exposure declarations; a selected
numerical output supplies its model and approximation identities; a causal
analysis accepts or refuses the intervention meaning; and a value object
retains the transformation and the complete uncertainty object. A simpler
baseline is a manually reviewed exported result with the same declarations.
Only a measured reduction in repeated translation errors or effort, with
equivalent scientific output, supports replacing that baseline with an
automated connection. No general MacroFinance--BayesFilter--causal adapter
has been qualified by the arithmetic example.

\section{Numerical authority and the error in a value transformation}

A connection to a state-space provider also needs a mathematical account of
approximation. Suppose, conditional on an explicitly declared model and
available history, the provider's target distribution over a finite state
$S$ is $P$ and its numerical approximation is $\widetilde P$. Let $f(S)$
be the value contribution, with known bounds $m\leq f(s)\leq M$.
Define total variation by
\[
 \operatorname{TV}(P,\widetilde P)
     =\frac12\sum_s|P(s)-\widetilde P(s)|.
\]
Then a supplied bound $\operatorname{TV}(P,\widetilde P)\leq\epsilon$
implies
\begin{equation}
 |E_Pf-E_{\widetilde P}f|\leq(M-m)\epsilon.
 \label{exec13:i3-tv}
\end{equation}
For a direct proof, let $\delta_s=P(s)-\widetilde P(s)$. The sum of the
positive $\delta_s$ and minus the sum of the negative $\delta_s$ both
equal total variation, since $\sum_s\delta_s=0$. Subtract the constant
$m$ from $f$, and bound the positive contribution by
$(M-m)\sum_{\delta_s>0}\delta_s$. Interchanging $P$ and $\widetilde P$
gives the other direction. This bound is a project derivation for the
specified finite problem, not a claim that an inspected filter supplies
$\epsilon$ automatically.

In the funding example, let the low and high states have coefficients
$0.6$ and $0.9$. Suppose the declared target probability of the high state
is $0.4$, while the approximate provider uses $0.5$. Their mean
coefficients are $0.72$ and $0.75$, and total variation is $0.1$.
The target conditional value is $30{,}260.75$ USD and the approximate
value is $33{,}188.28$ USD. Their difference $2{,}927.53$ equals the
bound in Equation~\eqref{exec13:i3-tv}, because the value range is
$29{,}275.31$. Exact rational coefficient arithmetic and a common-uniform
coupling check the calculation independently~\citep{exec13_i3_approximation}.
Neither distribution here is output from an executed BayesFilter run.

This construction shows what the consumer would need. If only the
approximate value and a justified $\epsilon=0.1$ bound were supplied,
the corresponding symmetric outer value interval would be
$[30{,}260.75,36{,}115.81]$ USD. That interval concerns numerical error
relative to the declared target distribution. It is not a frequentist
confidence interval for a bank effect and is not a posterior credible
interval for the random state. Model uncertainty, parameter uncertainty
and identification ambiguity require their own treatment. If the payoff is
unbounded, the displayed finite-range proof is unavailable; truncation,
tail control or another error metric would have to be justified rather
than assumed.

For an actual exchange, the causal-side mapping should therefore retain
the following fields in distinct roles. The model/data/coordinate identity
specifies which conditional distribution is being approximated. The
filter identity and approximation semantics identify the numerical
procedure. Its diagnostics describe what was observed in that procedure,
while an error statement must specify a norm, target and validity regime.
Finally, the value mapping gives the payoff, its horizon and any bounds
used to translate numerical error. A convergence trace is useful diagnostic
evidence but cannot silently become a total-variation certificate. When no
justified error bound is available, the consumer should display a
conditional approximate result and the missing error qualification.

A similar distinction applies to posterior output. A posterior distribution
computed under a structural model may support Bayesian decision analysis
conditional on that model and its priors. Passing posterior quantiles into
a field labeled frequentist coverage changes the asserted guarantee. Passing
a likelihood under a merely predictive model into an intervention-value
operation changes the scientific interpretation even earlier. The proposed
common contract should keep these tags and premises with the object, so a
consumer can compare a conditional structural calculation with a design-based
alternative without presenting them as interchangeable evidence.

This suggests a useful progression for the first numerical-provider trial.
Start with a finite model whose state law and value are exactly enumerable.
Check the declared approximation and the monetary mapping against that
reference. Next vary numerical settings while keeping the target fixed,
recording error, memory and time separately. Only then consider a richer
model with an explicit replacement for the finite reference. Such a sequence
tests the connection's mathematical meaning at low cost. It leaves the
intervention assumptions and the usefulness of its precision visible for a
later substantive decision.

\section{Servicing: preserve both policies and both decision times}

A servicing question asks how a feasible sequence of actions changes value.
It cannot be answered merely by importing an account's observed outcome or
a retrospective value gap. To make the required exchange concrete, consider
the following generated two-action law. The binary baseline covariate $X$
has probability $1/2$ of either value. Action $A_1$ occurs first; the binary
intermediate response $L$ then becomes available before action $A_2$; the
binary final recovery indicator $Y$ arrives last. Under interventions on
the two actions, independent exogenous uniforms generate
\begin{align}
 \Pr(L=1\mid X=x,A_1=a_1)&=0.2+0.2x+0.2a_1,\\
 \Pr(Y=1\mid x,a_1,l,a_2)&=0.1+0.2x+0.3l+0.2a_2+0.1a_1.
 \label{exec13:i3-servicing-law}
\end{align}
All stated probabilities lie between zero and one. Recovery earns 100 USD;
the first and second actions cost 6 and 4 USD respectively. For this short
illustration discounting is zero, and the total reward is
$R=100Y-6A_1-4A_2$. The independence and intervention interpretation are
definitions of this generated structural law, not conclusions from a fitted
observational regression.

Compare the policy $\pi_0$ taking neither action with the adaptive policy
\[
 \pi_1(x)=1-x,\qquad \pi_2(x,a_1,l)=1-l.
\]
The first action can use $X$ but not the later $L$. The second action may
use $L$ because its availability is part of the design. Write $p_L$ and
$p_Y$ for the two displayed probabilities. Direct integration gives
\begin{equation}
 V(\pi)=\frac12\sum_{x=0}^{1}\sum_{l=0}^{1}
   p_L(x,\pi_1(x))^l\{1-p_L(x,\pi_1(x))\}^{1-l}
   \left[100p_Y(x,\pi_1(x),l,\pi_2)-6\pi_1(x)-4\pi_2\right],
 \label{exec13:i3-policy-value}
\end{equation}
with $\pi_2$ evaluated at the displayed history. The sum integrates over
the intermediate response distribution induced by the first action. Holding
the observed distribution of $L$ fixed while changing $A_1$ would calculate
a different object. Substitution yields $V(\pi_0)=29$ and
$V(\pi)=43.6$, so the proposed contrast is $14.6$ USD per account.

An independent enumeration checks the same quantity through an observed
behavior law. Let $\Pr(A_1=1\mid X)=1/2$ and let
$\Pr(A_2=1\mid X,A_1,L)$ equal $1/4$ when $L=1$ and $3/4$ when
$L=0$. Generate actions with additional independent uniforms, so sequential
exchangeability and positivity hold by construction. If $e_1,e_2$ denote
the probabilities of the actions actually taken, then
\begin{equation}
 V(\pi)=E\left[
  \frac{\mathbf1\{A_1=\pi_1(X),\ A_2=\pi_2(X,A_1,L)\}}
       {e_1(A_1\mid X)e_2(A_2\mid X,A_1,L)}R\right].
 \label{exec13:i3-policy-ipw}
\end{equation}
To see the identity, factor the joint law of $(X,A_1,L,A_2,Y)$ and cancel
the two action probabilities on histories matching the policy. Summing over
$Y$ then gives its conditional mean, leaving
Equation~\eqref{exec13:i3-policy-value}. Exact rational enumeration of all
32 observed cells verifies total mass one, both values and the contrast~\citep{exec13_i3_examples}. This is a check of the generated finite law and
the proposed semantics. It supplies neither empirical coverage nor an
implemented longitudinal consumer adapter.

The corresponding FinMarketing exchange needs two feature snapshots or an
equivalent explicit history representation, two action definitions, the
availability of $L$, the cash horizon and the training identity of any fitted
policy. A policy chosen using the final evaluation outcomes is a further
target-selection problem, even when every field passes a schema. Candidate
policy learning should use an independent training sample or a separately
justified evaluation design. A filtering provider can assist in constructing
a declared observable history; it cannot make a privately observed state
available to the bank or remove unmeasured confounding by returning a latent
state estimate.

CardNPV illustrates a different target boundary. Its inspected
\texttt{CardTargetValue} uses historical value minus perfect-execution
value, supports the declared horizons of 24, 48, 120 and 240 months, and
distinguishes finite-only from conditional continuation~\citep{exec13_i3_inspection}.
Suppose, at a common admitted horizon, historical value is 900 USD and
perfect-execution value is 1,100 USD. That target is $-200$ USD. If a
particular feasible policy has value 980 USD, its contrast against history
is $80$ USD. Negating the accounting gap gives $200$, which still refers
to the perfect-execution intervention rather than the feasible policy.
An adapter must therefore compare the interventions, populations, horizons
and continuation assumptions before relating the targets. The generated
two-action example has its own short horizon; it is not supplied to CardNPV
as though it were an admitted 24-month account calculation.

Three alternatives remain useful. A scenario valuation can hold a supplied
policy response fixed and clearly label its assumptions. A longitudinal
causal analysis can attempt to identify the response from declared data and
assignment conditions. A structural state model can supply a conditional
intervention calculation under stronger dynamics. These alternatives answer
different evidential questions. The proposal should compare their assumptions,
information cost and diagnostics before selecting one for an actual consumer.
If the local service supports only a static target, the correct response to
this sequential request is an explicit unsupported-operation result together
with the preserved request; it must not return a static ATE under the same
policy label.
